\documentclass[11pt,a4paper]{article}

%% PACHETE MATEMATICE
\usepackage{amsmath,amssymb,amsfonts}
\usepackage{amsthm}
\usepackage{mathpazo}
\usepackage{eulervm}
\usepackage{enumerate}
\usepackage{color}
\usepackage{graphicx}
\usepackage[all]{xy}
\usepackage{xcolor}
\usepackage{framed}
\usepackage[unicode]{hyperref}
\setcounter{MaxMatrixCols}{20}
\usepackage{multicol}
%\usepackage[romanian]{babel}


%% ASPECT
\linespread{1.05}
\usepackage[T1]{fontenc}
\usepackage[utf8x]{inputenc}
\usepackage[margin=2cm]{geometry}
% \usepackage[color]{showkeys}
\usepackage{anyfontsize}
\usepackage{titlesec}
\titleformat*{\section}{\large\bfseries}

\usepackage{fancyhdr}
\pagestyle{fancy}
\rhead{\small \color{gray} Notițe de Adrian Manea}
\lhead{\small \color{gray} ETTI-AG, 2017---2018, A. Niță}


\usepackage[autostyle = false, style = german]{csquotes}
\usepackage[romanian]{babel}
\newcommand\qq{\enquote}

%% COMENZILE MELE
\renewcommand*{\proofname}{Dem.:}
\newcommand\bb{\mathbb}
\newcommand\dr{\mathrm}
\newcommand\kal{\mathcal}
\newcommand\fk{\mathfrak}
\newcommand\op{\oplus}
\newcommand\ot{\otimes}
\newcommand\ds{\displaystyle}
\newcommand\id{\indent}
\newcommand\nid{\noindent}
\newcommand\seq{\subseteq}
\newcommand\sneq{\subsetneq}
\newcommand\speq{\supseteq}
\newcommand\spneq{\supsetneq}
\newcommand\rar{\rightarrow}
\newcommand\Rar{\Rightarrow}
\newcommand\xrar{\xrightarrow}
\newcommand\lar{\leftarrow}
\newcommand\Lar{\Leftarrow}
\newcommand\xlar{\xleftarrow}
\newcommand\lrar{\leftrightarrow}
\newcommand\Lrar{\Leftrightarrow}
\newcommand\ideal{\trianglelefteq}
\newcommand\ai{\text{ a.\^{i}. }}
\newcommand\sk{(\textit{Schi\c{t}\u{a}}) }
\newcommand\rhup{\rightharpoonup}
\newcommand\rhdn{\rightharpoondown}
\newcommand\lhup{\leftharpoonup}
\newcommand\lhdn{\leftharpoondown}
\newcommand\vid{\emptyset}
\newcommand\wh{\widehat}
\newcommand\ol{\overline}
\newcommand\NN{\mathbb{N}}
\newcommand\RR{\mathbb{R}}
\newcommand\ZZ{\mathbb{Z}}
\newcommand\QQ{\mathbb{Q}}
\newcommand\CC{\mathbb{C}}

%% STILURILE MELE
\newtheoremstyle{dotless}{}{}{\itshape}{}{\bfseries}{:}{ }{}

\theoremstyle{dotless}
\newtheorem{theorem}{Teorem\u{a}}[section]
\newtheorem{proposition}{Propozi\c{t}ie}[section]
\newtheorem{lemma}{Lem\u{a}}[section]
\newtheorem{corollary}{Corolar}[section]
\newtheorem{exercise}{Exerci\c{t}iu}

\newtheoremstyle{dotlessdr}{}{}{}{}{\bfseries}{:}{ }{}
\theoremstyle{dotlessdr}
\newtheorem{example}{Exemplu}[section]
\newtheorem{definition}{Defini\c{t}ie}[section]
\newtheorem{remark}{Observa\c{t}ie}[section]




\begin{document}

\begin{center}
  {\large\textbf{Seminar suplimentar \\ Recapitulare pentru parțial}}
\end{center}

\vspace{1cm}

1. Să se determine coordonatele vectorului $u \in \RR^3$ în baza canonică,
dacă în baza
\[
  B = \{ (2, 1, 2), (1, 1, 0), (3, 0, 0) \}
\]
el are coordonatele $(4, 0, 5)$.

\vspace{1cm}

2. În spațiul real $\RR^3$, considerăm vectorii:
\[
  x_1 = (1, 0, 2), x_2 = (4, 1, -1), y_1 = (1, 2, 3), y_2 = (0, 5, 4), %
  y_3 = (1, 1, 1)
\]
și spațiile vectoriale:
\[
  V_1 = \mathrm{Sp}(x_1, x_2), \quad V_2 = \mathrm{Sp}(y_1, y_2, y_3).
\]

Să se determine o bază pentru $V_1, V_2, V_1 + V_2, V_1 \cap V_2$ și
un subspațiu $U \seq \RR^3$, astfel încît $U \oplus V_1 = \RR^3$.

\vspace{1cm}

3. Fie $V_1 = \{ P \in \RR[X]_3 \mid P(1) = 0 \}$ și subspațiul:
\[
  V_2 = \mathrm{Sp}(X + 8, X^2, 2X + 3X^2 + 4X^3).
\]

Să se determine cîte o bază în spațiile vectoriale reale $V_1, V_2$,
$V_1 + V_2, V_1 \cap V_2$ și un subspațiu $V \seq \RR[X]_3$ (dacă există),
astfel încît $V \oplus (V_1 + V_2) = \RR[X]_3$.

\vspace{1cm}

4. Fie aplicația liniară:
\[
  f : \RR^3 \to \RR^2, \quad f(x, y, z) = (x - 2y, z - 8y).
\]

Să se determine nucleul și imaginea sa și cîte o bază în fiecare subspațiu.
Completați baza nucleului pînă la o bază a lui $\RR^3$.

\vspace{1cm}

5. Fie aplicația liniară:
\[
  f : M_2(\RR) \to \RR[X]_2, \quad %
  f(\begin{pmatrix}
    a & b \\
    c & d
  \end{pmatrix}) = aX^2 + (b - c)X + 8d.
\]

\begin{enumerate}[(a)]
\item Găsiți matricea asociată aplicației liniare în bazele canonice;
\item Determinați $\mathrm{Ker}f$ și $\mathrm{Im}f$ și verificați teorema
  rang-defect.
\end{enumerate}

\vspace{1cm}

6. Fie matricea:
\[
  A = \begin{pmatrix}
    2 & 1 & 3 \\
    1 & 0 & a \\
    2 & 2 & 4
  \end{pmatrix}, \quad a \in \RR.
\]
\begin{enumerate}[(a)]
\item Găsiți aplicația liniară $f : \RR^3 \to \RR^3$, a cărei matrice
  în baza canonică este matricea $A$ de mai sus;
\item Găsiți $a \in \RR$ astfel încît aplicația determinată mai sus să fie
  izomorfism;
\item Determinați defectul aplicației $f$ (discuție după $a$).
\end{enumerate}

\vspace{1cm}

7. Fie $V = \RR[X]_2$ și mulțimea:
\[
  B' = \{ p_1 = 3X^2 + 2, p_2 = X - 5, p_3 = X^2 + 4 \}.
\]
\begin{enumerate}[(a)]
\item Arătați că $B'$ este bază a lui $V$;
\item Determinați matricea de trecere de la baza canonică la baza $B'$;
\item Fie $f : V \to \RR[X]_1$ morfismul de derivare (în raport cu $X$).
  Determinați matricea lui $f$ în baza canonică a lui $V$, față de baza canonică
  alui $\RR[X]_1$, precum și matricea lui $f$ în baza $B'$ a lui $V$
  față de baza $B^{\prime\prime} = \{ q_1 = 7, q_2 = X - 2 \}$ a lui $\RR[X]_1$.
\end{enumerate}

\vspace{1cm}

8. Fie spațiul $V = \RR^3$ și mulțimile:
\begin{align*}
  S_1 &= \{ (x, y, z) \in \RR^3 \mid x + y + z = 0 \} \\
  S_2 &= \{ (x, y, z) \in \RR^3 \mid x = 4z - y \text{ și } z = 2x - y \}
\end{align*}

\begin{enumerate}[(a)]
\item Arătați că $S_1 \hookrightarrow V$ și $S_2 \hookrightarrow V$;
\item Determinați $S_1 \cap S_2$ și o bază în acest subspațiu;
\item Este suma $S_1 + S_2$ directă?
\item Determinați $S_1 + S_2$ și o bază în acest subspațiu.
\end{enumerate}

\vspace{1cm}

9. Să se determine vectorii și valorile proprii pentru matricele:
\[
  A = \begin{pmatrix}
    1 & 0 \\
    8 & 2
  \end{pmatrix}, \quad
  B = \begin{pmatrix}
    0 & 2 & 2 \\
    2 & 0 & -2 \\
    2 & -2 & 0
  \end{pmatrix}, \quad
  C = \begin{pmatrix}
    0 & 1 & -2 \\
    1 & 1 & 1 \\
    1 & -1 & -1
  \end{pmatrix}.
\]

\vspace{1cm}

10. Fie aplicația liniară $f : M_2(\RR) \to M_2(\RR)$, definită prin:
\[
  f \begin{pmatrix}
    a & b \\
    c & d
  \end{pmatrix} = \begin{pmatrix}
    2a + b + d & 2b + 4c + 5d \\
    2c + d & 8d
  \end{pmatrix}
\]

Să se arate că $f$ nu este diagonalizabilă.

\vspace{1cm}

11. Determinați vectorii și valorile proprii pentru aplicațiile:
\begin{enumerate}[(a)]
\item $f : \RR[X]_3 \to \RR[X]_3, f(P) = 2P'$;
\item $f : \RR^3 \to \RR^3, f(a, b, c) = (a - b + c, b, b - c)$.
\end{enumerate}


\end{document}
